\documentclass{article}

\usepackage{amsmath}
\usepackage{amsthm}
\usepackage{amssymb}
\usepackage{tikz}
\usetikzlibrary{arrows.meta, positioning}
\theoremstyle{definition}
\newtheorem{definition}{Definíció}[section]

\title{Hőtan}
\author{Simon László}
\date{12/03/2026}

\begin{document}

\maketitle
\section{}
\begin{definition}[Hőállapot]
    Részecskék hőmozgása.
\end{definition}

\begin{definition}[Hőérzet]
    Szubjektív, befolyásolható
\end{definition}

\subsection{Hőmérséklet mérése}

\subsubsection{Fahrenheit}
\subsubsection{Celsius}
%Drawing missing
\begin{itemize}
    \item 100 fok C víz forráspontja
    \item 0 fok C víz fagyáspontja
\end{itemize}

\subsubsection{Hőmérők}
%Drawing missing
\begin{itemize}
    \item 212 fok F - 100 fok C - k.b. 373 fok K
    \item 32 fok F - 0 fok C - k.b. 273 fok K
\end{itemize}

\subsection{Hőtágulás}

%Kísérlet, ábra missing
\begin{itemize}
    \item Sín \begin{itemize}
        \item Melegben locsolják
        \item Közök hagyása - zakatolás
    \end{itemize}
    \item Villanyvezetékek - Betonsúlyok
    \item Csövek %insert tube picture here
    \item Hidak - Görgők %insert bridge picture here
\end{itemize}

Mitől függ?
\begin{equation}
    \Delta l = \alpha \cdot l_{0} \cdot \Delta T \\
\end{equation}
$\Delta l$ - hosszváltozás \\
$l_{0}$ - eredeti hossz \\
$\Delta T$ - hőmérsékletváltozás \\
$\alpha$ - anyagi minőség\\ \\
Hőtágulási együttható:
\begin{equation}
    [\alpha]=\frac{1}{C} = \frac{1}{K}
\end{equation}


%Note to self: draw function drawn in class with tikz


\end{document}